\documentclass[conference]{IEEEtran}

\IEEEoverridecommandlockouts
% The preceding line is only needed to identify funding in the first footnote. If that is unneeded, please comment it out.
\usepackage{cite}
\usepackage{amsmath,amssymb,amsfonts}
% \usepackage{algorithm}
\usepackage[ruled,vlined]{algorithm2e}
% \usepackage{algorithmic}
\usepackage{algpseudocode}
\usepackage{graphicx}
\usepackage{textcomp}
\usepackage{xcolor}
\usepackage{booktabs}
\usepackage{array}
\usepackage{makecell}
\usepackage{multirow}
\usepackage{subfig}


\usepackage{amsmath}

\def\BibTeX{{\rm B\kern-.05em{\sc i\kern-.025em b}\kern-.08em
    T\kern-.1667em\lower.7ex\hbox{E}\kern-.125emX}}
\begin{document}

\title{Improving the Reliability of 7/9-Rate Sparse Codes for STT-MRAM Using a Single-Label Neural Network\\}

% {\footnotesize \textsuperscript{*}Note: Sub-titles are not captured in Xplore and
% should not be used}
% \thanks{Identify applicable funding agency here. If none, delete this.}
% }

\author{
\IEEEauthorblockN{
Thanh-Loc Nguyen-Van\textsuperscript{1},
Nhon Nguyen-Quang\textsuperscript{1},
Huy Giang-Vinh\textsuperscript{1},
Tan Do-Duy\textsuperscript{1},
Chi Dinh-Nguyen\textsuperscript{2}
}

\IEEEauthorblockA{\textsuperscript{1}
Department of Computer and Communication Engineering\\
Ho Chi Minh City University of Technology and Education (HCMUTE)\\
Ho Chi Minh City, Vietnam
}

\IEEEauthorblockA{\textsuperscript{2}
Faculty of Computing Fundamentals\\
FPT University\\
Hanoi, Vietnam
}
}

\maketitle

\begin{abstract}

This paper 

\end{abstract}


\begin{IEEEkeywords}
Spin-torque transfer magnetic random access memory (STT-MRAM), sparse codes, neural network, BCH codes, Asymmetric write error.
\end{IEEEkeywords}



\section{Introduction}

Spin-Transfer Torque Magnetic Random Access Memory (STT-MRAM) is a promising non-volatile memory technology because of its high density, fast read/write operations, high endurance, and suitability for emerging Internet of Things (IoT) and artificial intelligence (AI) applications \cite{girard2020survey}. However, ensuring the reliability of STT-MRAM remains a major challenge. Both write-induced and read-induced errors can occur, and these errors become more severe as technology continues to scale \cite{hemaram2025asymmetric}.

The reliability of STT-MRAM is affected by several factors. Process variations in the magnetic tunnel junction (MTJ), including variations in its geometry and resistance, change the switching behavior and cause deviations in the switching current. In synchronous systems, the write time is fixed, so some memory cells may not complete the switching process within the allocated write period. In addition, thermal fluctuations can randomly change the MTJ state, while read operations may unintentionally disturb the stored data. Manufacturing imperfections further increase these effects. As device dimensions continue to shrink, these error sources become more severe, making reliability a major challenge for STT-MRAM \cite{wu2020survey}.


To improve STT-MRAM reliability, channel coding techniques are widely used to detect and correct memory errors. Existing approaches can be broadly classified into constrained coding, error-correcting codes (ECCs), and joint constrained coding with ECC \cite{nguyen2026error}. Constrained coding reduces error-prone data patterns, while ECC techniques, such as Hamming and BCH codes, improve reliability by correcting bit errors. However, conventional ECC schemes provide the same protection for both bit-flipping directions and therefore cannot fully exploit the asymmetric error characteristics of STT-MRAM. To address this limitation, sparse coding has been proposed to reduce the number of "1"s in the stored data, thereby reducing the number of $(0 \rightarrow 1)$ write transitions, which are more likely to fail than $(1 \rightarrow 0)$ transitions \cite{nguyen2021design}.

In \cite{nguyen2021design}, a 7/9-rate sparse code was proposed to improve the reliability of STT-MRAM by exploiting its asymmetric write-error characteristic. The proposed sparse code reduces the number of $(0 \rightarrow 1)$ write transitions, thereby lowering the write failure probability and improving the reliability of write operations. However, the scheme only reshapes the input data and does not provide error-correction capability. As a result, its performance is limited when multiple errors occur, especially under severe channel conditions.

Recently, machine learning-based decoding has shown promising performance in improving STT-MRAM reliability by learning the nonlinear characteristics of memory channels \cite{gruber2017deep}. Motivated by these advances, this paper extends the 7/9-rate sparse coding scheme proposed in \cite{nguyen2021design} by developing a Single-Label Neural Network (SLNN)-based sparse decoder and integrating sparse coding with ECC, as suggested in \cite{nguyen2026error}. In the proposed framework, sparse coding mitigates asymmetric write errors by reducing the number of $(0 \rightarrow 1)$ write transitions, while ECC corrects residual channel errors. As a result, the proposed framework enhances the reliability of STT-MRAM by effectively addressing both asymmetric write failures and read errors.






\begin{table*}[!t]
\centering
\caption{Comparison with existing neural-assisted reliability enhancement schemes for STT-MRAM.}
\label{tab:comparison}

\footnotesize
\setlength{\tabcolsep}{3pt}
\renewcommand{\arraystretch}{1.15}

\begin{tabular}{|l|c|p{2.2cm}|p{2.5cm}|p{2.0cm}|p{4.5cm}|}
\hline
\textbf{Paper} &
\textbf{Sparse Code} &
{\textbf{Sparse Decoder}} &
{\textbf{Neural Detector}} &
{\textbf{Error-Correction Capability}} &
{\textbf{Main Novelty}} \\
\hline \hline

Nguyen (2021) \cite{nguyen2021design} &
7/9-rate &
Maximum-likelihood &
$\times$ &
$\times$ &
Design of 7/9 sparse code. \\ \hline

Nguyen et al. (2021) \cite{nguyen2021performance} &
7/9-rate &
Maximum-likelihood &
RNN &
$\times$ &
RNN channel detector. \\ \hline

Nguyen et al. (2023) \cite{nguyen2023neural}&
7/9-rate &
Pearson &
MLP &
$\times$ &
MLP channel detector. \\ \hline

Nguyen (2026) \cite{nguyen2025stacking}&
4/9-rate &
Pearson &
Stacking ensemble &
Single-bit error correction &
Ensemble channel detector. \\ \hline

\textbf{This work} &
\textbf{7/9-rate} &
\textbf{Single-Label Neural Network (SLNN)} &
\textbf{$\times$} &
\textbf{Supports multi-bit error correction via BCH} &
\textbf{Neural sparse decoder integrated with BCH for enhanced STT-MRAM reliability.} \\

\hline 
\end{tabular}
\end{table*}





Table~\ref{tab:comparison} summarizes representative neural-assisted reliability enhancement techniques for STT-MRAM. Early work by Nguyen~\cite{nguyen2021design} proposed the 7/9-rate sparse code together with a conventional maximum-likelihood decoder, without exploiting neural networks. Subsequently, neural networks were introduced exclusively for channel detection. Nguyen \textit{et al.}~\cite{nguyen2021performance} employed the recurrent neural network (RNN)-based detector, while Nguyen \textit{et al.}~\cite{nguyen2023neural} adopted the multilayer perceptron (MLP) detector to improve sensing reliability before sparse decoding. More recently, an ensemble-learning detector was proposed for the 4/9-rate sparse code with single-bit error-correction capability~\cite{nguyen2025stacking}. Despite these advances, all existing studies utilize conventional sparse decoders after neural detection and therefore remain constrained by the intrinsic error-correction capability of the sparse code.

In contrast, the proposed method introduces neural learning into the sparse decoding stage itself. Specifically, an SLNN-based sparse decoder replaces conventional decoding algorithms and is integrated with an outer BCH decoder to provide multi-bit error-correction capability for STT-MRAM reliability enhancement. 





Therefore, the main contributions of this paper are summarized as follows:
\begin{itemize}
    \item First, we propose an SLNN-based decoder for the 7/9-rate sparse code by reformulating sparse decoding as a single-label multi-class classification problem. The proposed decoder directly maps noisy sparse codewords to their corresponding valid codewords, replacing conventional distance-based sparse decoding.
    \item Second, we develop an integrated BCH and sparse coding framework to enhance the reliability of STT-MRAM. In the proposed framework, the BCH code provides error-correction capability, while the sparse code mitigates asymmetric write errors by reducing the number of $(0 \rightarrow 1)$ write transitions.
    \item Finally, simulation results demonstrate that the proposed SLNN-based sparse decoder achieves an average BER improvement of approximately 63.3\% over the conventional maximum-likelihood (ML) sparse decoder across the considered operating range. Furthermore, the proposed integrated coding framework consistently achieves superior BER and FER performance compared with frameworks employing either the BCH code or the sparse code alone.

\end{itemize}

The remainder of this paper is organized as follows. Section~\ref{section_theory} introduces the STT-MRAM technology and the cascaded STT-MRAM channel model. Section~\ref{section_design} presents the overall system architecture and the proposed SLNN-based 7/9-rate sparse decoder. Section~\ref{section_simulation} provides the simulation results and corresponding discussions. Finally, Section~\ref{section_conclusion} concludes the paper.


\section{Channel model for STT-MRAM} \label{section_theory}
\subsection{STT-MRAM storage technology}

As illustrated in Fig.~\ref{figcell}, an STT-MRAM cell employs a conventional 1T--1MTJ architecture, consisting of one magnetic tunnel junction (MTJ) and one access nMOS transistor. The MTJ functions as the non-volatile storage element, while the access transistor controls both read and write operations \cite{cai2017cascaded}. The MTJ comprises two ferromagnetic layers separated by a thin insulating tunnel barrier. The reference layer (RL) has a fixed magnetization direction, whereas the free layer (FL) can switch its magnetization under the spin-transfer torque effect. The relative alignment of these two layers determines the resistance of the MTJ. When the RL and FL are aligned in the parallel (P) state, the MTJ exhibits a low resistance, denoted by $R_P$, which represents logic 0'. In contrast, when they are in the antiparallel (AP) state, the MTJ has a high resistance, denoted by $R_{AP}$, corresponding to logic 1'. Therefore, binary information is stored by distinguishing between these two resistance states \cite{nguyen2023improving}.

\begin{figure}[t]
\centerline{\includegraphics [scale=0.14]{fig/01.drawio.pdf}}
\caption{STT-MRAM structure. (a) STT-MRAM cell. (b) Parallel (P) state. (c) Antiparallel (AP) state \cite{nguyen2023improving}.}
\label{figcell}
\end{figure}



\begin{figure}[t]
\centerline{\includegraphics [scale=0.14]{fig/03.drawio.pdf}}
\caption{Writing process in STT-MRAM. (a) Bit-0 writing process. (b) Bit-1 writing process \cite{nguyen2021design}.}
\label{figwriteread}
\end{figure}

The write operation in STT-MRAM is inherently asymmetric, as illustrated in Fig.~\ref{figwriteread}. To write a logic `0', $V_{DD}$ is applied to the WL and BL while the source line (SL) is grounded, driving the MTJ to the low-resistance parallel (P) state. To write a logic `1', the current direction is reversed, switching the MTJ to the high-resistance antiparallel (AP) state. Because of the asymmetric switching behavior of the MTJ, the transition to the AP state generally requires more switching energy and a longer switching time than the transition to the P state. Consequently, the write error probability of bit-1 can be up to three orders of magnitude higher than that of bit-0, making write reliability a major challenge in STT-MRAM-based memory systems \cite{nguyen2021design}.

During the read operation, the word line (WL) is first activated, and a sensing current, much smaller than the write current, is passed through the MTJ to avoid disturbing the stored data. The resulting resistance is then compared with a reference threshold to determine the stored bit value. Read reliability mainly depends on the tunneling magnetoresistance (TMR) ratio, which measures the resistance difference between the P and AP states. A higher TMR ratio provides a larger sensing margin, thereby reducing read decision errors and improving the overall reliability of the memory system \cite{nguyen2021design}.




\subsection{Cascaded channel model}

To evaluate signal processing techniques for STT-MRAM, an accurate channel model is required to characterize the physical error mechanisms of the memory. Several channel models have been proposed in the literature \cite{cai2017cascaded,li2016efficient}. In this work, the cascaded STT-MRAM channel model introduced by Cai and Immink \cite{cai2017cascaded} is adopted.

\begin{figure}[!t]
\centering
\includegraphics[width=\columnwidth]{fig/04_v2.drawio.pdf}
\caption{Cascaded STT-MRAM channel model under the write-1 operation \cite{cai2017cascaded}.}
\label{fig3}
\end{figure}

Let $P_1$, $P_0$, and $P_r$ denote the write error probability for $1\rightarrow0$ transitions, the write error probability for $0\rightarrow1$ transitions, and the read disturb error rate, respectively. As shown in Fig.~\ref{fig3}, the channel consists of three cascaded components. A binary asymmetric channel (BAC) models asymmetric write errors with crossover probabilities $P_0/2$ and $P_1/2$. The read disturb error, which is unidirectional, is modeled by a Z channel. Finally, read decision errors are represented by a Gaussian mixture channel (GMC), where the MTJ resistances are modeled as Gaussian random variables: $R_0 \sim \mathcal{N}(\mu_0,\sigma_0)$ for the low-resistance state and $R_1 \sim \mathcal{N}(\mu_1,\sigma_1)$ for the high-resistance state.

For analysis, the BAC and Z channels are combined into an equivalent BAC-ZC channel with crossover probabilities

\begin{equation}
\begin{aligned}
p_0 &= \frac{P_0}{2} + \left(1 - \frac{P_0}{2}\right) P_r, \quad
q_0 = \left(1 - \frac{P_0}{2}\right)(1 - P_r), \\
p_1 &= \frac{P_1}{2} + (1 - P_r), \quad
q_1 = \left(1 - \frac{P_1}{2}\right) + \frac{P_1}{2} P_r.
\end{aligned}
\end{equation}

where $p_0$ and $p_1$ are the crossover probabilities for $0\rightarrow1$ and $1\rightarrow0$ transitions, respectively, while $q_0$ and $q_1$ denote the probabilities of correct transitions ($0\rightarrow0$ and $1\rightarrow1$).




\section{Proposed SLNN-Assisted Reliability Enhancement Framework} \label{section_design}
\subsection{System model}

Fig.~\ref{figsystem} illustrates the block diagram of the proposed system. The framework combines an outer BCH code, which provides error correction capability, with an inner 7/9-rate sparse code that shapes the input data to reduce write errors in STT-MRAM.

At the transmitter, the user data $\mathbf{u}$ is first encoded by the BCH encoder to generate the codeword $\mathbf{c}_1$. The BCH code provides error-correction capability against channel errors. The encoded sequence is then processed by the 7/9-rate sparse encoder to produce the sparse codeword $\mathbf{c}_2$. By reshaping the input data into sparse patterns, the sparse encoder reduces the number of high-energy write operations, thereby lowering the write error probability in STT-MRAM.

The sparse codeword is transmitted through the cascaded STT-MRAM channel, which models asymmetric write errors, read disturbance errors, and read decision errors. At the receiver, the received signal $\mathbf{r}$ is first decoded by the proposed SLNN-based sparse decoder. Unlike conventional hard-decision decoding, the proposed decoder exploits the soft information of the received signal to accurately recover the sparse codeword, producing the estimated BCH codeword $\hat{\mathbf{c}}_1$.

Finally, $\hat{\mathbf{c}}_1$ is passed to the BCH decoder, which corrects the remaining residual errors to recover the original user data $\hat{\mathbf{u}}$. By combining data shaping through the 7/9-rate sparse code, neural-assisted sparse decoding, and BCH error correction, the proposed framework effectively improves the reliability of STT-MRAM. For brevity, the detailed encoding and decoding procedures of the BCH and 7/9-rate sparse codes are omitted in this paper, as they have been comprehensively described in \cite{lin2001error} and \cite{nguyen2021design}, respectively.




\begin{figure*}[t]
\centerline{\includegraphics [scale=0.18]{fig/figsys_01.drawio.pdf}}
\caption{Block diagram of the proposed system.}
\label{figsystem}
\end{figure*}


\subsection{Single-Label Neural Network-Based 7/9-Rate Sparse Decoder}

For the $(7,9)$ sparse code introduced in \cite{nguyen2021design}, each 7-bit information word is uniquely mapped to a 9-bit codeword, yielding a codebook of 128 valid codewords, each having a Hamming weight of either 2 or 4. Accordingly, the decoding problem is formulated as a 128-class single-label classification task, where each codeword is assigned a unique class label. During training, the class labels are represented using one-hot encoding and used as the target outputs for supervised learning.

Fig.~\ref{fig:SLNN_decoder} illustrates the proposed single-label neural network (SLNN)-based decoder. The received signal vector $\mathbf{r}=[r_1,r_2,\ldots,r_9]$ is directly fed into a feedforward neural network (FFNN), which learns the nonlinear relationship between the noisy channel observations and the corresponding codeword classes. The output layer consists of 128 neurons, each representing one valid codeword in the sparse-code codebook. During inference, the decoder selects the class corresponding to the highest output probability as the predicted codeword. Finally, the predicted codeword is mapped to its corresponding 7-bit information word using a look-up table (LUT).


In this work, we adopt the feedforward neural network architecture proposed in \cite{gruber2017deep}. The network follows a $9$-$128$-$64$-$32$-$128$ configuration, consisting of three hidden layers with 128, 64, and 32 neurons, respectively. Each hidden layer comprises a fully connected layer followed by a LeakyReLU activation function, defined as
\begin{equation}
\phi(x)=
\left\{
\begin{array}{ll}
x, & x \ge 0,\\
\alpha x, & x < 0.
\end{array}
\right.
\label{eq:leakyrelu}
\end{equation}
where $\alpha=0.01$ is the negative slope coefficient. The output layer consists of 128 neurons, each corresponding to one valid codeword class. Let $\mathbf{z}=[z_0,z_1,\ldots,z_{127}]$ denote the output logits produced by the network. The predicted codeword class is determined by
\begin{equation}
\hat{c}=\arg\max_{0\leq i\leq127} z_i,
\label{eq:argmax}
\end{equation}
where $\hat{c}$ denotes the estimated class label corresponding to the decoded codeword.


Following the data generation strategy in \cite{bennatan2018deep}, the training dataset is constructed through Monte Carlo simulations of the STT-MRAM channel. Specifically, each valid 9-bit sparse codeword in the predefined codebook is transmitted through the channel to generate its corresponding noisy received vector. The resulting received vector is used as the input to the neural network, while the target label is defined as the index of the transmitted codeword in the codebook.

% The network input is given by
% %
% \begin{equation}
%     \mathbf{y}
%     =
%     [y_1,y_2,\ldots,y_9]
%     \in\mathbb{R}^{9},
%     \label{eq:received_vector}
% \end{equation}
% %
% where the nine elements are the read values obtained at the
% output of the STT-MRAM channel. The received vector is directly
% normalized without applying a hard-threshold decision:
% %
% \begin{equation}
%     \widetilde{\mathbf{y}}
%     =
%     \frac{\mathbf{y}-1.5}{0.5}.
%     \label{eq:input_normalization}
% \end{equation}

% The neural decoder employs a
% $9$-$128$-$64$-$32$-$128$ architecture, in which the
% $128$-$64$-$32$ hidden-layer configuration follows the neural
% decoder studied in \cite{gruber2017deep}. The three hidden
% layers perform the transformations $9\rightarrow128$,
% $128\rightarrow64$, and $64\rightarrow32$, respectively.
% Each linear layer is followed by batch normalization and a
% LeakyReLU activation function with a negative slope of $0.01$.
% The final layer performs the $32\rightarrow128$ transformation
% and produces 128 logits. No batch normalization or LeakyReLU
% activation is applied to this layer.

% Given the logit vector $\mathbf{z}\in\mathbb{R}^{128}$, the
% class probabilities are determined using the softmax function.
% The decoded class is selected as
% %
% \begin{equation}
%     \widehat{c}
%     =
%     \underset{c\in\{0,\ldots,127\}}
%     {\operatorname{arg\,max}}
%     \left[
%         \operatorname{softmax}(\mathbf{z})
%     \right]_{c}.
%     \label{eq:class_decision}
% \end{equation}
% %
% Finally, the predicted class index $\widehat{c}$ is converted
% into its corresponding 7-bit information word using the
% decoding lookup table. During training, each sample is assigned
% a single label, and the cross-entropy loss is computed directly
% from the 128 output logits.


% \begin{table}[tb]
% \centering
% \caption{Summary of the considered BCH–sparse concatenated coding schemes and their corresponding overall coding rates.}
% \label{tab4}
% \renewcommand{\arraystretch}{1.2}
% \begin{tabular}{|c|c|c|}
% \hline
% No. 
% & Coding scheme 
% & Overall rate 
%  \\
% \hline\hline

% 1 & BCH(127,106,3) + 7/9-rate sparse code
% & 0.649 
%  \\ \hline

% 2 & BCH(127,99,4) + 7/9-rate sparse code
% & 0.606
%  \\ \hline

% 3 & BCH(127,106,3) + 4/9-rate sparse code
% & 0.371
%  \\ \hline

% 4 & BCH(127,99,4) + 4/9-rate sparse code 
% & 0.346
%  \\
% \hline
% \end{tabular}
% \end{table}



\section{SIMULATION RESULTS AND DISCUSSION} \label{section_simulation}

In this section, the bit error rate (BER) and frame error rate (FER) performances of the proposed scheme are evaluated through Monte Carlo simulations. For each operating point, independent information frames were generated, encoded, transmitted through the STT-MRAM channel, and decoded using the corresponding decoding scheme. The BER was obtained by comparing the decoded information bits with the transmitted bits, while the FER was defined as the ratio of incorrectly decoded frames to the total number of transmitted frames. To ensure statistically reliable BER and FER estimates, each simulation was terminated after processing at least $10^6$ information bits.


\subsection{Simulation setup}

This subsection evaluates the BER performance of the considered coding schemes as functions of the coefficient of variation $\sigma_0/\mu_0$ and the write-error probability $P_1$, following the simulation methodology in \cite{an2025sparse,nguyen2021design}. The simulation parameters are adopted from the experimentally validated STT-MRAM model in \cite{cai2017cascaded,zhang2011stt}, using a 45-nm predictive technology model (PTM)-based STT-MRAM cell. The mean resistances of the low- and high-resistance states are set to $1~\mathrm{k}\Omega$ and $2~\mathrm{k}\Omega$, respectively, with identical coefficients of variation, i.e., $\sigma_0/\mu_0=\sigma_1/\mu_1$.

The parameter $\sigma_0/\mu_0$ characterizes resistance variation during read operations, where a larger value results in more severe read decision errors. In contrast, $P_1$ denotes the probability of a bit-1 write error caused by incomplete switching or read-disturb-induced state changes. Since $P_0$ and $P_r$ are typically much smaller than $P_1$ (i.e., $P_0,P_r\ll P_1$), the overall BER performance is mainly dominated by $P_1$.



\subsection{Simulation results}
\subsubsection{Bit Error Rate Performance}

In this subsection, we evaluate the BER performance of the proposed framework from two perspectives. First, the BER performance of the conventional 7/9-rate sparse code with maximum-likelihood (ML) decoding \cite{nguyen2021design} is compared with that of the proposed SLNN-based 7/9-rate sparse decoder. Second, the proposed integrated coding framework, which combines BCH coding with the proposed SLNN-based 7/9-rate sparse decoder, is compared with frameworks employing either BCH coding alone or the proposed SLNN-based 7/9-rate sparse code alone.

The simulation results are presented in Fig.~\ref{ber_vs_sigma} and Fig.~\ref{ber_vs_p1}, which illustrate the BER performance as a function of the coefficient of variation $(\sigma_0/\mu_0)$ and the write-error probability $(P_1)$, respectively. The results demonstrate that the proposed SLNN-based sparse decoder reduces the average BER by 63.3\% compared with the conventional ML-based sparse decoder over the considered operating range. Furthermore, the proposed integrated BCH and SLNN-based sparse coding framework consistently achieves superior BER performance compared with frameworks employing either BCH coding alone or the proposed SLNN-based sparse coding scheme alone.

\begin{figure}[t]
\centerline{\includegraphics [scale=0.65]{fig/BER_vs_sigma.png}}
\caption{BER performance versus the coefficient of variation $(\sigma_0/\mu_0)$.}
\label{ber_vs_sigma}
\end{figure}

\begin{figure}[t]
\centerline{\includegraphics [scale=0.65]{fig/BER_vs_p1.png}}
\caption{BER performance versus the write error rate $(P_1)$.}
\label{ber_vs_p1}
\end{figure}

\subsubsection{Frame Error Rate Performance}


In this subsection, we evaluate the frame error rate (FER) performance of the proposed coding framework. Similar to the BER evaluation, the conventional 7/9-rate sparse code with maximum-likelihood (ML) decoding is first compared with the proposed SLNN-based 7/9-rate sparse decoder. In addition, the proposed integrated coding framework, which combines BCH coding with the proposed SLNN-based sparse decoder, is compared with frameworks employing either BCH coding alone or the proposed SLNN-based sparse coding scheme alone.

The simulation results are presented in Fig.~\ref{fer_vs_sigma} and Fig.~\ref{fer_vs_p1}, which illustrate the FER performance as a function of the coefficient of variation $(\sigma_0/\mu_0)$ and the write-error probability $(P_1)$, respectively. As expected, the FER performance exhibits trends consistent with the BER results. Specifically, the proposed SLNN-based sparse decoder consistently outperforms the conventional ML-based sparse decoder over the entire range of operating conditions. Moreover, integrating BCH coding with the proposed SLNN-based sparse decoder provides the best FER performance among all considered schemes by effectively exploiting the complementary error-correction capabilities of the outer BCH code and the inner sparse coding scheme. These results demonstrate the effectiveness of the proposed integrated framework in improving frame-level reliability under both process variations and asymmetric write errors in STT-MRAM.



\begin{figure}[t]
\centerline{\includegraphics [scale=0.65]{fig/FER_vs_sigma.png}}
\caption{FER performance versus the coefficient of variation $(\sigma_0/\mu_0)$.}
\label{fer_vs_sigma}
\end{figure}


\begin{figure}[t]
\centerline{\includegraphics [scale=0.65]{fig/FER_vs_p1.png}}
\caption{BER performance versus the write error rate $(P_1)$.}
\label{fer_vs_p1}
\end{figure}

\section{Conclusion} \label{section_conclusion}

This paper 

\section*{Acknowledgment}







% -------------------------------------------------------------------------
\bibliographystyle{IEEEtran} % ieeetr
\bibliography{refs}





\vspace{12pt}
% \color{red}
% IEEE conference templates contain guidance text for composing and formatting conference papers. Please ensure that all template text is removed from your conference paper prior to submission to the conference. Failure to remove the template text from your paper may result in your paper not being published.

\end{document}
